%% kapa: exam=1 round=24 year=2013 period=1 form=A subject=economics subject_ko=경제원론 seq=18 number=58
%% source: Q-Net 감정평가사 기출문제 (cst00309) artlSeq=5073802 file=extracted/1008/5073802/2013년 제24회 감정평가사 1차시험 기출문제/1교시(A형).hwp
%% conversion: kordoc 4.24.51 -> tools/gen1.py (auto); review=proofread-sonnet; flags=none
\begin{question}{58}
두 재화 $X$와 $Y$를 소비하고 있는 甲의 효용함수는 $U(x, y)=\min\{4x+y, x+7y\}$이다. 수평축을 $X$의 소비량($x$), 수직축을 $Y$의 소비량($y$)이라고 할 때, 甲이 $X$ 재화 ４단위, $Y$ 재화 ７단위를 소비하는 점에서 무차별곡선의 기울기는?
\begin{choices}
  \item $-4$
  \item $- \frac{ 4 } { 7 }$
  \item $- \frac{ 1 } { 4 }$
  \item $- \frac{ 1 } { 7 }$
  \item $- \frac{ 7 } { 4 }$
\end{choices}

eviewflag{question 58 not rendered in rhwp pages; not verified against image}
\end{question}
\answer{1}
